\input{../KU-Beamer-Template/style/template.tex}
\usepackage{../KU-Beamer-Template/style/koc} 
\usepackage{minted}
\usepackage{upquote}
\usepackage{graphicx}

% pdflatex --shell-escape lecture5.tex & pdflatex --shell-escape lecture5.tex

\title{KOLT Python}
% https://www.geeksforgeeks.org/difference-method-function-python/
\subtitle{Functions} 
\newdate{date}{21}{10}{2019}
\date{\displaydate{date}}
\author{Ahmet Uysal}

\titlegraphic{\includegraphics[scale=0.2]{../KU-Beamer-Template/style/images/logo_kolt.pdf}}

\setbeamercovered{invisible} % transparent
\makeatletter
\let\@@magyar@captionfix\relax
\makeatother
\begin{document}
    \maketitle

    \frame{\frametitle{Agenda}\tableofcontents}

    \section{Recap}
       \begin{frame}{While Loops}
            Repeat some \texttt{<expression>}s \underline{as long as} a \texttt{<condition>} is \texttt{True}.
            \begin{columns}
                \begin{column}{0.45\textwidth}
                    \inputminted[frame=single,framesep=2pt]{python3}{../Lecture3/code-examples/while1.py}
                \end{column}
                \begin{column}{0.45\textwidth}
                    \inputminted[frame=single,framesep=2pt]{python3}{../Lecture3/code-examples/while2.py}
                \end{column} 
            \end{columns}
            \LARGE
            \inputminted[frame=single,framesep=2pt]{python3}{../Lecture3/code-examples/while3.py}
            \texttt{<condition>} is only checked \textbf{\underline{before}} each execution.
        \end{frame}

        \begin{frame}{List Slicing}
            \Large
            Access collection of elements with \textbf{\texttt{[start:stop:step]}}\\
            Gives a list, even when number of elements is not bigger than 1.
            \inputminted[frame=single,framesep=2pt,firstline=10]{python3}{../Review1/code-examples/slicing.py}                                
            \LARGE
            Slices with \texttt{step} = 1 are called \textbf{Basic Slice}.\\
            Slices with \texttt{step != 1} are called \textbf{Extended Slice}.
        \end{frame}

        \begin{frame}{len() Function}
          \pause
          \LARGE
          \texttt{len()} is an operator to determine the size of lists, strings, etc.
          \pause
          \bigskip
          \inputminted[frame=single,framesep=2pt]{python3}{../Lecture3/code-examples/length.py}
        \end{frame}

        \begin{frame}{range() Function}
            \pause
            \LARGE
            \texttt{range(\textbf{start, stop, step})} is a function to create ranges
            \bigskip
            \inputminted[frame=single,framesep=2pt]{python3}{../Lecture3/code-examples/range.py}
        \end{frame}
      
        \begin{frame}{For Loops}
            \pause
            \begin{columns}
                \begin{column}{0.5\textwidth}
                    \inputminted[frame=single,framesep=2pt]{python3}{../Lecture3/code-examples/for1.py}
                \end{column}
               \pause 
                \begin{column}{0.5\textwidth}
                    \inputminted[frame=single,framesep=2pt]{python3}{../Lecture3/code-examples/for2.py}
                    \pause
                    \inputminted[frame=single,framesep=2pt]{python3}{../Lecture3/code-examples/for3.py}
                    \pause
                    \inputminted[frame=single,framesep=2pt]{python3}{../Lecture3/code-examples/for4.py}
                \end{column} 
            \end{columns}
        \end{frame}
        
        \begin{frame}{Break, Continue \& Pass}
            \begin{columns}
                \begin{column}{0.5\textwidth}
                    \textbf{\texttt{break}} immidiately terminates the closest loop
                    \bigskip  
                    \inputminted[frame=single,framesep=2pt]{python3}{../Lecture3/code-examples/break1.py}
                    \pause
                    \inputminted[frame=single,framesep=2pt]{python3}{../Lecture3/code-examples/break2.py}
                \end{column}
               \pause 
                \begin{column}{0.5\textwidth}
                    \textbf{\texttt{continue}} skips to the next iteration of the loop
                    \bigskip  
                    \inputminted[frame=single,framesep=2pt]{python3}{../Lecture3/code-examples/continue1.py}
                    \pause
                    \inputminted[frame=single,framesep=2pt]{python3}{../Lecture3/code-examples/continue2.py}
                \end{column} 
            \end{columns}
        \end{frame}
        
        \begin{frame}{Break, Continue \& Pass}
            \pause
            \LARGE
            \textbf{\texttt{pass}} does not have an effect
            \bigskip  
            \inputminted[frame=single,framesep=2pt]{python3}{../Lecture3/code-examples/pass.py}
            \pause
            \begin{itemize}
                \item Loops, conditional statements, functions etc. cannot be empty 
            \end{itemize}
        \end{frame}


    \section{Functions}
        \subsection{Defining Functions}
        \begin{frame}{Functions}
            \pause
            \LARGE
            Functions are blocks of \textbf{organized}, \textbf{reusable} code that carry some \textbf{specific} tasks.
            \pause
            \begin{itemize}
                \item \textbf{\texttt{input([prompt])}:} \\
                \pause
                If the prompt \textbf<5->{argument} is present, it is written to standard output without a trailing newline. The function then reads a line from input, converts it to a string (stripping a trailing newline), and \textbf<5->{returns} that. When \texttt{EOF} is read, \texttt{EOFError} is \textbf<5->{raised}.   
            \end{itemize}

            \pause
            \pause

            \begin{center}
                \tikzstyle{int}=[draw, fill=blue!20, minimum size=2em]

                \begin{tikzpicture}[node distance=2.5cm,auto]
                    \node [int] (a) {$function$};
                    \node (b) [left of=a,node distance=4cm, coordinate] {a};
                    \node [coordinate] (end) [right of=a, node distance=4cm]{};
                    \path[->] (b) edge node {$*arguments$} (a);
                    \draw[->] (a) edge node {\texttt{return}} (end) ;
                \end{tikzpicture}
            \end{center}
        \end{frame}   
 
        \begin{frame}{Defining Functions}
            \pause
            \LARGE
            \textbf{\texttt{def}} keyword introduces a function \textit{definition}.
            \normalsize
            \pause
            \inputminted[frame=single,framesep=2pt, lastline=8]{python3}{code-examples/function_def.py}
            \pause
            \inputminted[frame=single,framesep=2pt, lastline=8]{python3}{code-examples/function_def2.py}
            \pause
            \inputminted[frame=single,framesep=2pt, lastline=8]{python3}{code-examples/function_def3.py}
        \end{frame}


        \begin{frame}{Functions}
            \inputminted[frame=single,framesep=2pt]{python3}{code-examples/function_ex.py}
            \pause
            \inputminted[frame=single,framesep=2pt]{python3}{code-examples/function_ex2.py}
        \end{frame}

        \begin{frame}{Functions}
            \LARGE
            \textit{Defining} a \texttt{function} only makes it available.\\
            \pause 
            You should \textit{call} the \texttt{function} to execute.\\           
            \pause
            \newline
            \textbf{\texttt{fib\_100 = fibonacci\_series(100)}}\\
            \pause
            \textbf{\texttt{what\_is\_going\_on = print(fib\_100)}}\\
            \newline
            \pause
            \textbf{Every} function returns \textbf{one} value! \pause Functions implicitly return \texttt{None} if they complete without a return statement.
        \end{frame}

      \subsection{return Statement}

        \begin{frame}{Return}
            \LARGE
            \pause
            \inputminted[frame=single,framesep=2pt]{python3}{code-examples/return.py}
            \pause
            Return \textbf{immidiately} terminates the function. So, the output is 8.
        \end{frame}

      \subsection{Parameters}    

        \begin{frame}{Default Parameters}
            The values of parameters can be set to used as default.
            \newline
            In \textbf{\texttt{print(*args, sep=\textquotesingle \ \textquotesingle, end=\textquotesingle \textbackslash n\textquotesingle )}}, 
            sep and end has default values.
            \pause
            \inputminted[frame=single,framesep=2pt, lastline=15]{python3}{code-examples/default.py}
            \textbf{Valid Uses}
            \begin{columns}
                \begin{column}{0.50\textwidth}
                    \pause
                    \inputminted[frame=single,framesep=2pt, lastline=15]{python3}{code-examples/valid1.py}  
                \end{column}
                \begin{column}{0.50\textwidth}
                    \pause
                    \inputminted[frame=single,framesep=2pt, lastline=15]{python3}{code-examples/valid1_1.py}                    
                \end{column}
            \end{columns}
        \end{frame}

        \begin{frame}{Default Parameters}
            \inputminted[frame=single,framesep=2pt, lastline=15]{python3}{code-examples/default.py}
            \pause
            \LARGE
            \textbf{Invalid Usages}
            \inputminted[frame=single,framesep=2pt, lastline=15]{python3}{code-examples/valid2.py}  
        \end{frame}

        \begin{frame}{Variadic Positional Arguments}
            \pause
            \LARGE
            How to allow function to accept arbitrary number of arguments.
            \newline
            In \textbf{\texttt{print(*args, sep=\textquotesingle \ \textquotesingle, end=\textquotesingle \textbackslash n\textquotesingle )}}, 
            you can put as many args as you want.
            \pause
            \newline 
            \newline Suppose we want a \texttt{max} function that works as so:
            \newline max(3, 5) gives 5.
            \newline max(3, 4, 2) gives 4.
            \newline product(3, 5, -1, 2, 10, 20, 13, 34) gives 34.
            \newline 
        \end{frame}

        \begin{frame}{Variadic Positional Arguments: my\_max}
          \pause
          \large
          \inputminted[frame=single,framesep=2pt]{python3}{code-examples/variadic.py}
        \end{frame}
      
      \subsection{Variable Scope}

        \begin{frame}{Local \& Global Variables}
          \pause
          \large
          \begin{columns}
            \begin{column}{0.50\textwidth}
              \begin{itemize}
                \item Local variables are created in functions.
                
                \item Global variables are created out of the functions.
              \end{itemize}
              \inputminted[frame=single,framesep=2pt, lastline=15]{python3}{code-examples/var.py}
            \end{column}
            \begin{column}{0.50\textwidth}
              \pause
              \inputminted[frame=single,framesep=2pt, lastline=15]{python3}{code-examples/var2.py}                    
              \pause
              \inputminted[frame=single,framesep=2pt, lastline=15]{python3}{code-examples/var3.py}                    
            \end{column}
          \end{columns}
        \end{frame}
\end{document}